\chapter{Come calcola il solutore}
\label{ch:moc}

Il resto di questo documento parla di continuo di ``celle'',
``colonne'', ``piedi'', ``caratteristiche'' e ``compatibilità''.
Questo capitolo le definisce tutte partendo da zero. Descrive il
\emph{metodo delle caratteristiche}, il metodo numerico su cui posa
l'intero programma, e vale le quattro pagine che occupa: quasi ogni
lezione tecnica registrata più avanti --- compreso ogni difetto in
appendice --- è un'affermazione sugli oggetti qui introdotti.

\section{L'unico fatto fisico che il metodo sfrutta}

In un gas che si muove \emph{più lentamente} del suono, una
perturbazione si propaga in ogni direzione, a monte compreso:
l'intero campo di moto è accoppiato a sé stesso, e risolverlo
significa risolvere ovunque simultaneamente. In un gas che si muove
\emph{più velocemente} del suono, questo non può avvenire. Una
perturbazione prodotta in un punto viene trascinata a valle più in
fretta di quanto riesca a propagarsi, quindi la sua influenza è
confinata a un cuneo che si apre a valle di quel punto.

Quel cuneo ha un semiangolo preciso, l'\emph{angolo di Mach}
$\alpha$, dato da
\begin{equation}
\sin\alpha = \frac{1}{M}, \qquad M = \frac{q}{c},
\label{eq:machangle}
\end{equation}
dove $q$ è la velocità del flusso, $c$ la velocità locale del suono
e $M$ il loro rapporto, il \emph{numero di Mach}. Un flusso più
veloce dà un cuneo più stretto. I suoi due lati, tracciati per il
punto, sono le \emph{caratteristiche}: uno a un angolo
$\theta + \alpha$ rispetto all'asse e uno a $\theta - \alpha$, dove
$\theta$ è la direzione in cui il gas sta effettivamente scorrendo
(Figura~\ref{fig:mocprimer}a). Per convenzione la prima famiglia si
chiama $C^{+}$ e la seconda $C^{-}$.

Le caratteristiche non sono solo geometria. Lungo ciascuna di esse
il flusso deve soddisfare una \emph{relazione di compatibilità}:
un'equazione differenziale ordinaria che lega la variazione di
velocità alla variazione di direzione del flusso, con un termine
aggiuntivo nel flusso assialsimmetrico che tiene conto del condotto
che si allarga al crescere del raggio. Questo è l'intero contenuto
fisico che il solutore usa. Due famiglie di linee, una relazione
lungo ciascuna --- nessuna griglia di volume, nessun avanzamento nel
tempo, nessuna iterazione fino allo stato stazionario.

\section{Il processo elementare: una cella}

Il metodo trasforma quelle relazioni in aritmetica con una singola
mossa ripetuta, la \emph{cella} (o \emph{processo elementare}),
mostrata in Figura~\ref{fig:mocprimer}b.

Si supponga che due punti vicini siano già risolti: posizione e
velocità note in entrambi. Dal punto inferiore si traccia la sua
caratteristica $C^{+}$; da quello superiore la sua $C^{-}$. In
flusso supersonico queste due linee convergono e si incrociano in
esattamente un punto a valle. Quell'incrocio fornisce due equazioni,
una per linea, per le due coordinate del nuovo punto. Le due
relazioni di compatibilità, integrate lungo i due brevi segmenti
appena tracciati, danno altre due equazioni per le due componenti di
velocità lì. Quattro equazioni, quattro incognite $(x, y, u, v)$.

Le equazioni sono non lineari --- le direzioni caratteristiche
dipendono proprio dallo stato che si sta cercando --- quindi ogni
cella è chiusa da una piccola iterazione di Newton. Due dettagli su
come vengono chiuse contano più avanti:

\emph{Coefficienti al punto medio.} I coefficienti delle relazioni
di compatibilità sono valutati non all'inizio del segmento ma nel
suo \emph{punto medio}, usando la media dei due stati agli estremi.
È la regola dei trapezi; fa sì che l'errore per cella si riduca con
il quadrato del passo anziché linearmente, ed è il motivo per cui
l'intera marcia si dice del secondo ordine. È anche la ragione per
cui una cella non deve mai attraversare una regione in cui lo stato
cambia bruscamente: mediare i due estremi di un tale segmento è
allora privo di significato (il Capitolo~\ref{ch:plug} contiene un
esempio misurato).

\emph{Certificazione di Newton.} Una volta che l'iterazione di
Newton di una cella è convergita, il programma verifica che il
residuo delle quattro equazioni sia davvero sceso al livello
dell'arrotondamento in virgola mobile, scalato con la dimensione
delle incognite. Una cella che si arena viene colta e segnalata
invece di inquinare in silenzio la marcia. Ogni affermazione
``tutte le celle certificate da Newton'' in questo documento è
quella verifica, eseguita su ogni cella.

\section{Colonne, frontiere e piedi}

I nuovi punti crescono da punti risolti, quindi una marcia si
organizza in \emph{colonne}: da una colonna risolta, le celle
interne producono la colonna successiva, e la spazzata si ripete
verso valle (Figura~\ref{fig:mocprimer}c). Una marcia ha dunque
bisogno di un punto da cui cominciare --- una \emph{linea iniziale}
(o linea dei valori iniziali) attraverso il flusso, sulla quale lo
stato completo è fornito come dato. In un ugello di razzo quella
linea sta appena a valle della gola.

Ai bordi del flusso la cella cambia forma, perché una delle sue due
caratteristiche genitrici manca e una condizione al contorno ne
prende il posto. In questo documento compaiono tre frontiere.

\emph{Un asse di simmetria}: il flusso non può attraversarlo, il che
fornisce la relazione mancante.

\emph{Una parete assegnata}: la posizione del nuovo punto è già
nota, perché deve trovarsi su un contorno dato con una pendenza
data. La cella allora lavora \emph{all'indietro}. Deve trovare il
\emph{piede} --- il punto sulla colonna precedente da cui la
caratteristica in arrivo è partita --- che in generale cade fra due
punti memorizzati e va localizzato cercando lungo la colonna e
interpolando. La relazione di compatibilità viene poi chiusa lungo
la corda da quel piede al punto di parete. È questa la cella che
permette a un progettista di consegnare una parete arbitraria
(Capitolo~\ref{ch:wallinput}), e l'accuratezza della ricerca del
piede è un tema ricorrente.

\emph{Una frontiera di getto libero}: il caso in cui non c'è parete
alcuna e lo scarico è delimitato dall'atmosfera. Qui la situazione è
esattamente rovesciata: la \emph{pressione} è nota (è $p_a$, poiché
sul getto non spinge altro che l'atmosfera) mentre la
\emph{posizione} della frontiera è incognita e va risolta. È questo
l'unico pezzo di fisica genuinamente nuovo di cui l'ugello a spina
ha bisogno, e il Capitolo~\ref{ch:plug} lo costruisce e lo
certifica.

\begin{figure}[htbp]
\centering
\includegraphics[width=\textwidth]{fig_moc_primer.pdf}
\caption{Come calcola la marcia, per un lettore nuovo al metodo
delle caratteristiche. Gli assi sono ovunque posizione assiale $x$
[m] e raggio $y$ [m]. (a)~Per ogni punto di un flusso supersonico
passano due caratteristiche, inclinate dell'angolo di Mach $\alpha$
da una parte e dall'altra della direzione locale del flusso
$\theta$; lungo ciascuna, una relazione di compatibilità lega la
variazione di velocità alla variazione di angolo. (b)~La cella
interna: da due punti già risolti, l'incrocio di una $C^{+}$ e di
una $C^{-}$ fissa la posizione del nuovo punto, e le due relazioni
di compatibilità --- chiuse con coefficienti presi nei punti medi
dei segmenti ($\diamond$) --- ne fissano la velocità; un'iterazione
di Newton risolve il sistema $4\times4$ che ne risulta. (c)~Le celle
costruiscono colonne a partire da colonne. A una parete assegnata la
posizione è nota anziché calcolata, quindi la cella deve localizzare
per interpolazione il \emph{piede} della caratteristica in arrivo
sulla colonna precedente. Ogni lezione strutturale del
Capitolo~\ref{ch:plug} e dell'Appendice~\ref{ch:appendix} è
un'affermazione su queste tre figure.}
\label{fig:mocprimer}
\end{figure}

\section{Spigoli vivi: l'unica cosa che una cella non sa fare}

Dove una parete gira uno spigolo vivo, il flusso si espande
attraverso un \emph{ventaglio di espansione centrato}: una
distribuzione continua di caratteristiche tutte uscenti da quel
singolo punto geometrico, ciascuna portatrice di uno stato diverso.
Nello spigolo stesso il flusso è dunque a più valori e i gradienti
sono infiniti.

Nessuna cella ordinaria può scavalcare un tale ventaglio, per
entrambe le ragioni dette sopra: un segmento che attraversasse il
ventaglio medierebbe stati che differiscono per l'intera rotazione,
e vicino allo spigolo l'incrocio di caratteristiche che fissa la
nuova posizione degenera perché tutte le linee si incontrano in un
punto solo. Il rimedio è che un ventaglio centrato non ha bisogno di
essere marciato --- è \emph{autosimile}, quindi lo stato su ogni
raggio è noto in forma chiusa a partire dal solo angolo del raggio.
Lo spigolo viene perciò \emph{costruito} come dato e le celle
ordinarie riprendono fra i raggi, dove il flusso è di nuovo
regolare. Il Capitolo~\ref{ch:plug} mostra che cosa accade quando
questa regola viene violata, con il numero misurato.

\section{Gas reale, e da dove vengono le derivate}

Vale la pena enunciare qui, una volta per tutte, due scelte
implementative.

Il gas non è assunto caloricamente perfetto. Le proprietà
termodinamiche provengono da dati tabulati, e lo stato è recuperato
dalla velocità attraverso le relazioni di energia ed entropia.
Questo conta in un modo non ovvio che il
Capitolo~\ref{ch:wallinput} documenta: le proprietà tabulate sono
interpolate a tratti, quindi la derivata \emph{esatta} del codice
porta una piccola increspatura alla scala delle celle della tabella
--- una proprietà reale dell'obiettivo discreto, non un difetto, ma
tale da rendere mal posto un ingenuo test sulle derivate.

E l'intera marcia --- ogni cella, ogni iterazione di Newton, ogni
ricerca del piede --- è scritta in modo da poter essere
differenziata automaticamente. È questo che trasforma un solutore
nel motore di un ottimizzatore, ed è il soggetto di tutto ciò che
segue.
